\subsection{Purpose}

The \textbf{AgriHome Vision Console} is a web-based monitoring and decision-support product for controlled-environment agriculture and laboratory plant studies. Operators and researchers use it to observe tray-level imagery, track individual plants, review automated health assessments, organize trays into logical \textbf{meshes}, manage image-capture schedules, and operate \textbf{Raspberry~Pi / Klipper} edge capture devices. This \textbf{Architectural Design Specification (ADS)} records how the product is decomposed into layers and subsystems, how data moves among them, and which interfaces connect internal components to each other and to external partners.

This document is a companion to the team's \textbf{System Requirements Specification (SRS)} for the same project.\cite{srs} The SRS states \emph{what} the system must do for customers (for example, plant health presentation, automatic capture, environmental data, dashboards, authentication, and secure storage). This ADS states \emph{how} those responsibilities are allocated to software structure: presentation, application orchestration, processing (vision adapters, edge disease hook, schedule runner), persistent data, and hardware edge capture.

\subsection{Scope}

\textbf{In scope} for this ADS are: the logical five-layer architecture (presentation, application, processing, data, hardware); subsystem names, assumptions, responsibilities, and interface summaries; major data flows from user or edge-agent action to storage; identification of optional processing and search backends; Raspberry~Pi / Klipper edge integration at the architectural level; and deployment context at the level of application, database, object storage, and bench devices.\cite{edgeops}

\textbf{Out of scope} are: user-interface visual design mockups beyond structural responsibilities; exhaustive REST field specifications (see engineering docs); source file paths and build commands; formal threat models (see edge security notes); and training procedures for ML models.

\subsection{System summary}

AgriHome addresses limitations of ad hoc spreadsheets, disconnected camera feeds, and manual record-keeping by providing a \textbf{single console} that ties together:

\begin{itemize}
  \item \textbf{Identity and access} — authenticated users interact only with their permitted trays, plants, and meshes; edge devices authenticate with hashed device API keys.
  \item \textbf{Observation} — latest captures, charts, and monitoring history support plant status display, historical data, and dashboard-style summaries.
  \item \textbf{Edge capture} — Raspberry~Pi devices with Agri-Home/klipper (\texttt{fswebcam} via \texttt{save\_image.sh}) auto-provision, heartbeat, and deliver imagery via Take Picture (primary agent capture; optional HTTP streamer still) and scheduled multi-pose capture.\cite{klipper}
  \item \textbf{Interpretation} — health reports and optional model-backed inference (including post-ingest disease analysis) enrich raw imagery where configured.
  \item \textbf{Planning} — schedules express when captures should occur, including destination \texttt{raspberry-pi-edge}.
  \item \textbf{Persistence} — relational storage holds authoritative domain and edge state; object storage holds images referenced from that state.
\end{itemize}

The reference implementation uses a modern full-stack web platform (browser and installable PWA client, server-hosted application, PostgreSQL, file storage, Firebase Authentication, optional HTTP-based inference services, and Python \texttt{agrihome\_agent} on the Pi). The architecture in this document is described so that those technologies could be substituted behind the same layer boundaries where a future maintainer requires it.

\subsection{Document organization}

Section~2 presents the \textbf{project overview} and describes the five primary layers. Section~3 gives a \textbf{system description} with a consolidated subsystem-and-data-flow view. Sections~4--8 document \textbf{data}, \textbf{presentation}, \textbf{application}, \textbf{processing}, and \textbf{hardware} layer subsystems respectively, each with assumptions, responsibilities, and interface tables. The \textbf{References} section lists external citations and the companion SRS.
